% ============================================================================
%  chapter8.tex
%  Filesystems in Depth: Metadata, Links, and Directories
% ============================================================================

\chapter{Filesystems in Depth}

\section{The Filesystem Is a Namespace}

Chapter~5 treated files as ``a sequence of bytes you can read or
write,'' and that was the right first cut. Now we go one layer deeper:
a \emph{filesystem} is the kernel's system for \emph{naming,
organizing, and persisting} those byte sequences. Understanding it
means answering three questions the earlier chapter deferred: where do
the bytes physically live, what is a \emph{name}, and what does a
\emph{directory} really contain?

The three answers, in one picture:

\begin{verbatim}
   A filesystem has three layers:

     NAMES        a tree of directory entries:   /home/you/data.txt
        |
        v
     INODES       per-file metadata + a map of where the data is:
                  { type, permissions, size, [block 4][block 9]... }
        |
        v
     BLOCKS       fixed-size chunks of disk (or memory) that hold
                  the bytes themselves.

   - A NAME is a directory entry: a string in a directory file.
   - An INODE ("index node") is the kernel's record for ONE file:
     its type, its permissions, its size, and pointers to its blocks.
   - A BLOCK is the actual storage, typically 4 KB.
\end{verbatim}

The term to keep: a filesystem is a \emph{namespace}. The tree of
paths you type is the interface; behind it sits a machinery of inodes
and blocks that the kernel maintains and you almost never see. When we
say ``the file system,'' we mean this whole apparatus---not just the
directory tree.

\begin{osconcept}
Persistence and naming are an OS problem for the same reason every
service so far has been: only the kernel can give programs a
\emph{shared, protected, persistent} namespace. A library cannot make a
name survive a reboot; the filesystem does, because the kernel owns
the storage devices and the mapping from names to blocks. Providing
that namespace---and keeping it consistent across crashes and
processes---is one of the oldest responsibilities of an operating
system.
\end{osconcept}

\section{stat: Reading the Inode's Record}

The kernel exposes the contents of an inode through the \code{stat}
system call. It fills a \code{struct stat} with every fact about the
file that is not its data: its type, its permissions, its size, how
many hard links point to it, who owns it, and when it was touched.

\begin{lstlisting}[caption={stat: dump the metadata of a path.}]
#include <stdio.h>
#include <stdlib.h>
#include <sys/stat.h>
#include <sys/sysmacros.h>

int main(int argc, char *argv[]) {
    struct stat st;
    int r;

    if (argc != 2) {
        fprintf(stderr, "usage: %s path\n", argv[0]);
        return 1;
    }

    r = stat(argv[1], &st);
    if (r < 0) {
        perror("stat");
        return 1;
    }

    printf("size        : %ld bytes\n",  (long) st.st_size);
    printf("mode        : 0%o (octal)\n", (unsigned) st.st_mode);
    printf("hard links  : %ld\n",          (long) st.st_nlink);
    printf("owner       : uid=%d gid=%d\n", (int) st.st_uid, (int) st.st_gid);
    printf("device      : major=%d minor=%d\n",
           major(st.st_dev), minor(st.st_dev));
    printf("blocks      : %ld (x512 bytes)\n", (long) st.st_blocks);

    return 0;
}
\end{lstlisting}

\begin{verbatim}
   What each field is (the inode's record, in plain words):

     st_mode    a bit field packing BOTH the file type and the
                permissions. Two kinds of info in one integer.
     st_nlink   how many directory entries name this same inode.
                A file with no name left has nlink == 0 and is
                actually deleted (or waiting for its last close).
     st_size    the length in bytes of the DATA.
     st_uid, st_gid   who owns it (the numbers behind ls's names).
     st_mtime   last time the DATA changed.  ("m" = modified)
     st_ctime   last time the INODE metadata changed. (NOT creation)
     st_atime   last time it was read.
     st_blocks  how many 512-byte units of disk it occupies. Compare
                st_blocks*512 to st_size: a sparse file (ch.5) shows
                blocks*512 < size, because holes cost nothing.

   st_mode's two halves:

     bits 12-15 : the TYPE      bits 0-11 : the PERMISSIONS
     0100000 regular (S_IFREG)   0755     = rwx r-x r-x
     0040000 directory (S_IFDIR)
     0120000 symbolic link (S_IFLNK)
\end{verbatim}

The type bits are read through the macros \code{S\_ISREG},
\code{S\_ISDIR}, \code{S\_ISLNK}, and friends---never by comparing raw
numbers:

\begin{lstlisting}[caption={The S\_IS* macros decode the type half of st\_mode.}]
#include <stdio.h>
#include <sys/stat.h>

const char *type_name(mode_t mode) {
    if (S_ISREG(mode))  return "regular file";
    if (S_ISDIR(mode))  return "directory";
    if (S_ISLNK(mode))  return "symbolic link";
    if (S_ISCHR(mode))  return "character device";
    if (S_ISBLK(mode))  return "block device";
    if (S_ISFIFO(mode)) return "fifo";
    if (S_ISSOCK(mode)) return "socket";
    return "unknown";
}
\end{lstlisting}

One warning before we go on: there are \emph{two} stat-like calls.
\code{stat(path, ...)} follows a symbolic link and reports the
\emph{target's} inode. \code{lstat(path, ...)} reports the link itself.
The difference matters the moment we walk directories, because a
directory may contain links we must not chase.

\begin{osconcept}
An inode is kernel state: one per file, owned by the kernel, updated
by the kernel (permissions, timestamps, link counts) as the world
changes around you. \code{stat} is a read-only window into that state.
Only the kernel can maintain the mapping from names to inodes to
blocks, so ``what is this file?'' is answered by the OS, not by your
program.
\end{osconcept}

\section{File Types}

The type half of \code{st\_mode} distinguishes seven kinds of object,
and every one of them is ``a file'' in the eyes of the namespace.
This is the design decision of Chapter~5 paying off: a uniform
interface over wildly different objects.

\begin{verbatim}
   the seven types, as the first character of `ls -l`:

     -  regular file     bytes of data (the kind you usually mean)
     d  directory        a file whose bytes ARE directory entries
     l  symbolic link    a file whose bytes are a target path
     c  character device a device that transfers bytes one at a time
                          (terminals, /dev/null)
     b  block device     a device that transfers in blocks (disks)
     p  fifo             a named pipe from Chapter 7
     s  socket           a network endpoint from Chapter 17

   `ls -l /dev/null` shows:  c...
   `ls -ld /home`     shows:  d...
   `ls -l /proc`      shows:  d...  (and inside it, virtual files)
\end{verbatim}

\begin{verbatim}
   The surprising consequence:

     A directory is a FILE. Its "data" is a list of names -> inode
     numbers. `read()` a directory directly is awkward and
     discouraged, which is why we will use the opendir/readdir API
     (Section 8.5). But underneath, a directory is just bytes the
     kernel happens to interpret as entries.
\end{verbatim}

\begin{osconcept}
Treating devices, pipes, and directories as file types is the
\emph{everything is a file} doctrine, and it is an OS design decision
of enormous power. Because every object in the namespace answers the
same \code{stat}/\code{open}/\code{read} interface, tools never need
to know what they are touching. The kernel's job is to translate the
generic calls into whatever the underlying object actually needs---
hardware for a device, bytes for a regular file.
\end{osconcept}

\section{Permissions}

The low 12 bits of \code{st\_mode} describe who may do what. There are
three classes---user, group, other---and three operations per class---
read, write, execute. In octal, each class is one digit: 4 for read,
2 for write, 1 for execute, summed.

\begin{verbatim}
   mode 0644, drawn bit by bit:

       octal:     0     6       4       4
             special  u w x   g r w x   o r w x   (class: rwx)
       bits:  sticky  | | |    | | |    | | |
                     user     group    other
       0644 = rw- r-- r--

   The three special bits (not shown here) are setuid, setgid, and
   the sticky bit -- we mention them only to note they exist.

   mode 0755 (a program): rwx r-x r-x   -- owner may edit, all may run
   mode 0600 (a secret):  rw- --- ---   -- only the owner
   mode 0777 (a warning): rwx rwx rwx   -- everyone everything
\end{verbatim}

The kernel enforces these on every \code{open}: it compares the
requester's uid/gid against the file's owner and group, picks the
matching class, and checks the appropriate bits. This check happens
\emph{in the kernel}, at \code{open} time, once---which is why a file
can be opened and then its permissions changed without breaking the
already-open descriptor.

\begin{verbatim}
   Where the bits are consulted, and who changes them:

     open("f", O_RDONLY)   -> kernel checks read bits for your class
     open("f", O_WRONLY)   -> kernel checks write bits
     chmod("f", 0600)      -> you change the bits (if you own it)
     umask(022)            -> your process asks the kernel to mask
                              off bits from every new file's mode

   The umask explains a common mystery: you open() with mode 0666 and
   get 0644. The kernel ANDed your mode with ~umask (022) at open
   time. The kernel, not the library, owns this masking.
\end{verbatim}

\begin{lstlisting}[caption={Decoding permissions bit by bit, like ls does.}]
#include <stdio.h>
#include <sys/stat.h>

char file_type_char(mode_t mode) {
    if (S_ISREG(mode))  return '-';
    if (S_ISDIR(mode))  return 'd';
    if (S_ISLNK(mode))  return 'l';
    if (S_ISCHR(mode))  return 'c';
    if (S_ISBLK(mode))  return 'b';
    if (S_ISFIFO(mode)) return 'p';
    if (S_ISSOCK(mode)) return 's';
    return '?';
}

/* Print the 10-character ls -l permissions field. */
void print_mode_bits(mode_t mode) {
    const char *rwx = "rwxrwxrwx";
    int i;

    putchar(file_type_char(mode));

    /* Bit 8 = user read ... bit 0 = other execute. */
    for (i = 0; i < 9; i++) {
        int bit;

        bit = 8 - i;
        if (mode & (1 << bit)) {
            putchar(rwx[i]);
        } else {
            putchar('-');
        }
    }
}
\end{lstlisting}

\begin{osconcept}
The permission model is enforced by the kernel because users cannot be
trusted to enforce it on each other, and because the check must happen
at the exact moment of every \code{open}. The kernel compares real and
effective ids, walks the three classes, and refuses---returning
\code{EACCES}---before any data moves. Making permission checks part of
the syscall boundary is what makes multi-user sharing safe at all.
\end{osconcept}

\section{Links: Two Kinds of Indirection}

The \code{st\_nlink} field in our \code{stat} dump is a count of
\emph{hard links}. Understanding the two kinds of link is the key to
the inode model, and both kinds come from a single idea: a \emph{name}
and a \emph{file} are different things, and several names can point at
one file.

\subsection{Hard Links}

A hard link is a second name for the \emph{same inode}. Creating one
increments \code{st\_nlink}; deleting one decrements it. The data is
freed only when the count reaches zero. The two names are perfectly
equal---there is no ``real'' file and ``link.'' When you create a file,
\code{st\_nlink} is 1; \code{link()} makes it 2.

\begin{lstlisting}[caption={Creating both kinds of link.}]
#include <stdio.h>
#include <unistd.h>

int main(void) {
    if (link("data.txt", "data_hard.txt") < 0) {
        perror("link");
    }
    if (symlink("data.txt", "data_soft.txt") < 0) {
        perror("symlink");
    }
    return 0;
}
\end{lstlisting}

\begin{verbatim}
   HARD LINK: two directory entries, ONE inode.

   directory                   the single inode (nlink = 2)
   +--------------------+       +------------------------+
   | data.txt ----------+------>| type, perms, size,     |
   | data_hard.txt -----+       | and pointers to the    |
   +--------------------+       | actual data blocks     |
                                +------------------------+

   - Both names are the SAME file. Renaming one leaves the other.
   - Deleting data.txt only decrements nlink to 1; data survives.
   - Hard links CANNOT cross filesystems (they are inode numbers,
     and inode numbers are only meaningful within one filesystem),
     and cannot point at a directory.
\end{verbatim}

\subsection{Symbolic Links}

A symbolic link (symlink) is a different beast: a tiny file whose
\emph{content is a path string}. It is a pointer in the namespace, not
a second name for an inode. The kernel resolves the path whenever you
open through it. That is why a symlink \emph{can} cross filesystems,
\emph{can} point at a directory, and \emph{dangles} when its target
vanishes.

\begin{verbatim}
   SYMBOLIC LINK: a name containing a path.

   directory                       inodes
   +--------------------+       +------------------------+
   | data.txt ----------+------>| inode for data.txt     |
   | data_soft.txt -----+       +------------------------+
   +--------------------+             ^
        "data_soft.txt" is a file whose bytes are the STRING
        "data.txt". Opening it makes the kernel open data.txt.

   - stat()  follows the link (reports the target).
   - lstat() reports the link itself.
   - Rename or delete data.txt, and data_soft.txt still "exists"
     but resolves to nothing: a dangling link.
\end{verbatim}

\begin{osconcept}
Links are kernel bookkeeping about the namespace. A hard link is a
reference count on an inode; a symlink is a small file containing a
path. Only the kernel can maintain the inode table, resolve paths,
and decide when a file's data may finally be reclaimed. The \emph{nlink}
count in \code{stat} is literally a kernel-owned counter that your
program can only read.
\end{osconcept}

\section{Directories and Traversal}

A directory is a file whose contents are \emph{directory entries}: a
list of (name, inode number) pairs. To read it you use the directory
stream API: \code{opendir} opens a handle, \code{readdir} returns one
entry at a time, \code{closedir} releases it.

\begin{lstlisting}[caption={Listing a directory, entry by entry.}]
#include <stdio.h>
#include <stdlib.h>
#include <dirent.h>

int list_dir(const char *path) {
    DIR *d;
    struct dirent *e;

    d = opendir(path);
    if (d == NULL) {
        perror("opendir");
        return -1;
    }

    while (1) {
        e = readdir(d);
        if (e == NULL) {
            break;
        }
        if (e->d_name[0] == '.') {
            continue;    /* skip "." and ".." (and dotfiles here) */
        }
        printf("%s\n", e->d_name);
    }

    closedir(d);
    return 0;
}
\end{lstlisting}

\begin{verbatim}
   What readdir() hands you, one struct per entry:

     struct dirent {
         char d_name[256];   the entry's name
         ...                 (the inode number, and a type hint that
                              MAY be present -- never rely on it)
     };

   - Every directory contains "." (itself) and ".." (its parent);
     real listings skip them.
   - readdir gives you only NAMES. To get size/permissions you must
     stat() each one -- that is exactly why `ls -l` reads the whole
     directory AND stats every entry.
   - The order of entries is not meaningful and not sorted.
\end{verbatim}

Recursing into subdirectories turns a lister into a tree walker, and
that is where \code{lstat} becomes essential: using \code{stat} here
would follow symlinks and could loop forever on a link that points at
an ancestor.

\begin{lstlisting}[caption={walk.c: a recursive tree walker, using lstat.}]
#include <stdio.h>
#include <stdlib.h>
#include <string.h>
#include <dirent.h>
#include <sys/stat.h>

/* Recursively print every entry under base, indented by depth. */
void walk(const char *base, int depth) {
    DIR *d;
    struct dirent *e;
    int i;

    d = opendir(base);
    if (d == NULL) {
        return;
    }

    while (1) {
        struct stat st;
        char path[1024];
        int plen;

        e = readdir(d);
        if (e == NULL) {
            break;
        }
        if (strcmp(e->d_name, ".") == 0 || strcmp(e->d_name, "..") == 0) {
            continue;
        }

        plen = snprintf(path, sizeof(path), "%s/%s", base, e->d_name);
        if (plen < 0 || plen >= (int) sizeof(path)) {
            continue;          /* name too long; skip this entry */
        }

        /* lstat, not stat: a symlink must be reported AS a link. */
        if (lstat(path, &st) < 0) {
            continue;
        }

        for (i = 0; i < depth; i++) {
            printf("  ");
        }
        printf("%s\n", e->d_name);

        if (S_ISDIR(st.st_mode)) {
            walk(path, depth + 1);
        }
    }

    closedir(d);
}

int main(int argc, char *argv[]) {
    const char *root;

    root = (argc >= 2) ? argv[1] : ".";
    walk(root, 0);
    return 0;
}
\end{lstlisting}

\begin{verbatim}
   Output shape for a small tree:

     .
       src
         main.c
         util.c
       README.md
       build.sh

   Every level is an opendir/readdir/closedir cycle; every entry is
   lstat()ed. That is precisely the loop `find` and `du` run.
\end{verbatim}

\begin{osconcept}
Directory traversal is kernel-managed iteration. The directory is a
file whose format the kernel owns, so user code reads it only through
the kernel's \code{readdir} interface---never by guessing at bytes.
The \code{lstat}-vs-\code{stat} distinction exists because the kernel
must honor your choice of whether to follow a link, and getting that
choice wrong (as \code{stat} does in a walker) is how infinite
recursion and path loops enter a program.
\end{osconcept}

\section{/proc: The Virtual Filesystem}

\textbf{Concept:} a \emph{virtual filesystem} is a filesystem whose
``files'' are not stored anywhere---their contents are \emph{generated
on the fly} by the kernel when you read them. Linux's \code{/proc} is
the canonical example. Every running process has a directory
\code{/proc/<pid>} full of files describing it, manufactured from
kernel data structures the moment you open them.

\begin{verbatim}
   /proc is not a disk. It is a window into the kernel:

     /proc/self/status     the calling process's state (parent,
                           threads, memory, ...)
     /proc/self/cmdline    the command line that started it
     /proc/self/fd/        a directory of its open descriptors
     /proc/<pid>/fd/       the same, for any other process
     /proc/meminfo         whole-system memory numbers
     /proc/cpuinfo         what the CPUs are
     /proc/net/tcp         the kernel's TCP connection table

   Reading them is just read(): the bytes are composed on demand.
   They answer the OS questions of Chapters 4-6 from inside a C
   program -- without signals, pipes, or privilege tricks.
\end{verbatim}

\begin{lstlisting}[caption={Reading the kernel through a virtual file.}]
#include <stdio.h>
#include <stdlib.h>
#include <fcntl.h>
#include <unistd.h>

int main(void) {
    int fd;
    char buf[512];
    ssize_t got;

    fd = open("/proc/self/status", O_RDONLY);
    if (fd < 0) {
        perror("open");
        return 1;
    }

    got = read(fd, buf, sizeof(buf) - 1);
    if (got < 0) {
        perror("read");
        return 1;
    }
    buf[got] = '\0';
    printf("%s\n", buf);

    close(fd);
    return 0;
}
\end{lstlisting}

The output is ordinary text (its first lines name the pid and parent,
the thread count, and the memory sizes). Reading it is nothing but the
Chapter~5 \code{open}/\code{read}/\code{close} pattern applied to a
file that the kernel creates from its own tables---which is why
\code{/proc} is a wonderful teaching tool and a beloved debugging
instrument for systems programmers.

\begin{osconcept}
\code{/proc} shows the kernel \emph{reusing the file abstraction for
introspection}. The same namespace machinery, inodes, and read path
now serve as a read-only API to kernel state. This is an OS design
choice: rather than inventing a new interface for ``ask the kernel
about itself,'' Linux made the kernel a filesystem. It only works
because the file abstraction is general enough to serve both disks and
data structures.
\end{osconcept}

\begin{researchnote}[The paper that made filesystems fast]
The metadata and layout story of this chapter has a famous ancestor.
In 1984, Marshall Kirk McKusick and colleagues at Berkeley published
``A Fast File System for UNIX,'' describing the redesign of the
original Unix filesystem. Its innovations are precisely the ideas we
have been handling all chapter: separating inode metadata from data
blocks, grouping related blocks for sequential reads, and laying out
data to minimize disk seeks.

\begin{verbatim}
   The paper's central design principle (condensed):

       "The file system is organized to keep the disk head moving
        in the right direction: data for one file is placed in
        contiguous block groups, and the inode for a file sits near
        its data."

   In our chapter's terms: st_blocks, inode-vs-data separation, and
   the reason `ls -l` can be fast are all consequences of this layout
   design. The 4 KB block size used throughout this book descends
   directly from this paper's block-group scheme.
\end{verbatim}

Why read it: it is the clearest demonstration that a filesystem is a
\emph{performance-and-durability} design as much as a namespace, and
that the questions we still answer today (block size, metadata
placement, crash consistency) were first posed in this paper.

\paperref{M.~K.~McKusick, W.~N.~Joy, S.~J.~Leffler, R.~S.~Fabry, ``A Fast
File System for UNIX,'' ACM Transactions on Computer Systems 2(3),
1984}{https://www.cs.utexas.edu/~dahlin/classes/cs375/papers/fast-ffs.pdf}
\end{researchnote}

\section{Practical: A Mini ls -l}

Let us combine \code{opendir}, \code{lstat}, and our permission
decoder into a genuine tool: a small \code{ls -l} that lists a
directory with type, permissions, link count, size, and name---the
output you have read a thousand times, now understood byte by byte.

\begin{lstlisting}[caption={minils.c: a real directory listing tool.}]
#include <stdio.h>
#include <stdlib.h>
#include <dirent.h>
#include <sys/stat.h>

char file_type_char(mode_t mode) {
    if (S_ISREG(mode))  return '-';
    if (S_ISDIR(mode))  return 'd';
    if (S_ISLNK(mode))  return 'l';
    if (S_ISCHR(mode))  return 'c';
    if (S_ISBLK(mode))  return 'b';
    if (S_ISFIFO(mode)) return 'p';
    if (S_ISSOCK(mode)) return 's';
    return '?';
}

void print_permissions(mode_t mode) {
    const char *rwx = "rwxrwxrwx";
    int i;

    putchar(file_type_char(mode));
    for (i = 0; i < 9; i++) {
        int bit;

        bit = 8 - i;
        if (mode & (1 << bit)) {
            putchar(rwx[i]);
        } else {
            putchar('-');
        }
    }
}

int main(int argc, char *argv[]) {
    const char *path;
    DIR *d;
    struct dirent *e;

    path = (argc >= 2) ? argv[1] : ".";

    d = opendir(path);
    if (d == NULL) {
        perror("opendir");
        return 1;
    }

    while (1) {
        struct stat st;
        char full[1024];
        int plen;

        e = readdir(d);
        if (e == NULL) {
            break;
        }
        if (e->d_name[0] == '.') {
            continue;
        }

        plen = snprintf(full, sizeof(full), "%s/%s", path, e->d_name);
        if (plen < 0 || plen >= (int) sizeof(full)) {
            continue;
        }

        if (lstat(full, &st) < 0) {
            continue;
        }

        print_permissions(st.st_mode);
        printf(" %2ld %8ld %s\n",
               (long) st.st_nlink,
               (long) st.st_size,
               e->d_name);
    }

    closedir(d);
    return 0;
}
\end{lstlisting}

\begin{verbatim}
   $ ./minils src
   -rw-r--r--  1     142 main.c
   -rw-r--r--  1      89 util.c

   You are now reading those ten characters and three columns the way
   the kernel writes them: type bits, permission bits, a reference
   count, and a size from the inode -- fetched by lstat per entry.
\end{verbatim}

\section{Structures in Bytes}

The inode record that \code{stat} reads is a byte structure. On
x86-64 glibc, \code{struct stat} is 144 bytes; the fields this chapter
used sit at these offsets:

\begin{verbatim}
   struct stat (x86-64 glibc), the fields you met:

   offset  size  field        meaning
   0       8     st_dev       device the file lives on
   8       8     st_ino       inode number
   16      8     st_nlink     hard-link count
   24      4     st_mode      type + permission bits
   28      4     st_uid       owner
   32      4     st_gid       group
   40      8     st_size      length in bytes
   72      8     st_blocks    x512-byte disk units
   96      16    st_mtime     data last modified
   112     16    st_atime     last access
   128     16    st_ctime     metadata changed (NOT creation)
\end{verbatim}

The on-disk ext4 inode is 256 bytes and holds the same ideas in
compact form:

\begin{verbatim}
   ext4 inode, key offsets:

   offset  size  field
   0       2     mode
   2       2     uid
   4       4     size_lo
   8       4     atime
   12      4     ctime
   16      4     mtime
   24      2     gid
   26      2     links_count      (st_nlink's source)
   28      4     blocks_lo
   40      60    block[15]        the map to data blocks
   106     2     size_high        (size above 4 GB)
\end{verbatim}

And a directory entry is bytes too: 4 bytes of inode number, 2 of name
length, 1 type byte, then the name. The kernel parses these records;
\code{readdir} hands you the pieces, and \code{ls -l} reads the rest
from the 144-byte \code{stat} record above.

\section{Summary}

\begin{verbatim}
   The filesystem in one paragraph:

   names (directory entries) -> inodes (metadata + block map) -> blocks

   Tools learned this chapter:
     stat()/lstat()    read an inode's record
     S_IS*()           decode the type half of st_mode
     mode bits         the permission model (rwx, umask, chmod)
     link()/symlink()  hard vs symbolic indirection
     opendir/readdir   iterate a directory
     /proc             the kernel as a virtual filesystem

   Rules:
     1. lstat when walking directories, stat when you mean the target;
     2. readdir gives names only -- stat each entry for details;
     3. nlink==0 means the file's data can be reclaimed;
     4. permissions are enforced at open() time by the kernel.
\end{verbatim}

\begin{osconcept}
Everything in this chapter was the OS maintaining a namespace: inodes,
reference counts, permission checks, virtual files, and directory
iteration are all kernel-managed state behind the paths you type. The
filesystem is where the OS's oldest promise lives---``your data will
survive, named, and protected''---and understanding its internals is
what lets you build databases, editors, and tools that treat storage
as a real resource instead of magic.
\end{osconcept}

We now have files and the metadata around them. In the next chapter we
turn back to processes for their other half: how they are grouped and
supervised---process groups, sessions, and the discipline of writing
daemons that run forever without a terminal.

